\subsection{根系与由反射生成的群}

命题 8. 设 F 是一个 l 维实 Hilbert 空间。令 H 为 F 的仿射超平面集，G 为由正交反射 \(s_{H}\)（关于超平面 \(H \in H\)）生成的群。假设第五章第 3 节的条件成立（即都有 \(g(H) \in H\)（对所有 \(H \in H\) 和 \(g \in G\)），且 G 在 F 上适当作用）。还假设 0 是 G 的一个特殊点，且属于 G 的平移群 T 的秩为 l。则在 \(V = F^{*}\) 中存在唯一的约化根系 R，使得从 F 到 \(V^{*}\) 的典范同构将 G 变为 R 的仿射韦伊群 \(W_{a}\)。

首先注意，关于 T 的假设蕴含 G 是本质的：否则仿射空间 F 将分解为乘积 \(F_{0} \times F_{1}\)，其中 \(\dim F_{1} < l\)，而群 G 可与在 \(F_{1}\) 上适当作用的一个位移群相识别（第五章第 3 节第 8 号命题 6），于是 T 的秩不会是 l。

令 \(\mathfrak{H}_0\) 为 \(\mathbf{H} \in \mathfrak{H}\) 中满足 \(0 \in \mathbf{H}\) 的元素所成的集合。对于 \(\mathbf{H} \in \mathfrak{H}_0\)，令 \(\mathfrak{H}_{\mathbf{H}}\) 为 \(\mathfrak{H}\) 中平行于 \(\mathbf{H}\) 的元素所成的集合。由于 0 是特殊点，\(\mathfrak{H}\) 是各 \(\mathfrak{H}_{\mathbf{H}}\)（其中 \(\mathbf{H} \in \mathfrak{H}_0\)）的并。由于 \(\mathbf{T}\) 的秩为 \(l\)，存在 \(v \in \mathbf{F}\)，使得沿向量 \(v\) 的平移属于 \(\mathbf{T}\) 且 \(v \notin \mathbf{H}\)。超平面 \(\mathbf{H} + kv\)（当 \(k \in \mathbf{Z}\) 时）两两不同且属于 \(\mathfrak{H}_{\mathbf{H}}\)。现在令 \(a\) 为 \(F\) 中与 \(\mathbf{H}\) 正交的单位向量；则 \(\mathbf{H} + (v|a)a \in \mathfrak{H}_{\mathbf{H}}\)，并且由于 \(\mathfrak{H}\) 是局部有限的（第五章第 1 节第 1 号引理 1），存在最小的实数 \(\lambda > 0\)，使得 \(\mathbf{H} + \lambda a \in \mathfrak{H}\)。我们将证明，\(\mathfrak{H}_{\mathbf{H}}\) 正是超平面 \(\mathbf{H} + k\lambda a\)（当 \(k \in \mathbf{Z}\) 时）所成的集合。事实上，

\[
\mathrm{H} ^ {\prime} = \mathrm{H} + \lambda a \in \mathfrak {H} _ {\mathrm{H}}
\]

而元素 \(s_{\mathrm{H}'} \circ s_{\mathrm{H}}\) 属于 \(\mathbf{G}\)，是沿向量 \(2\lambda a\) 的平移（第五章第 2 节第 4 号命题 5）。因此，\(\mathrm{H} + 2n\lambda a = (s_{\mathrm{H}'}s_{\mathrm{H}})^n(\mathrm{H})\) 和 \(\mathrm{H} + (2n + 1)\lambda a = (s_{\mathrm{H}'}s_{\mathrm{H}})^n(\mathrm{H}')\) 属于 \(\mathfrak{H}_{\mathrm{H}}\)。另一方面，若 \(\mathrm{L} \in \mathfrak{H}_{\mathrm{H}}\)，则存在 \(\xi \in \mathbf{R}\) 使得 \(\mathrm{L} = \mathrm{H} + \xi \lambda a\)，并且存在整数 \(n\) 使得

\[
\text { either } 2 n <   \xi \leqslant 2 n + 1, \text { or } 2 n - 1 <   \xi \leqslant 2 n.
\]

在第一种情形，\((s_{\mathrm{H}}s_{\mathrm{H}^{\prime}})^{n}(\mathrm{L}) = \mathrm{H} + (\xi -2n)\lambda a\)，其中

\[
0 <   (\xi - 2 n) \lambda \leqslant \lambda
\]

且由 \(\lambda\) 的定义可知 \(\xi = 2n + 1\)；在第二种情形，

\[
s _ {\mathrm{H}} (s _ {\mathrm{H}} s _ {\mathrm {H^ {\prime}}}) ^ {n} (\mathrm{L}) = \mathrm{H} + (2 n - \xi) \lambda a \text {with} 0 \leqslant (2 n - \xi) \lambda <   \lambda
\]

且由 \(\lambda\) 的定义可知 \(\xi = 2n\)。

由此，若 \(\alpha_{\mathrm{H}}\) 是 \(F\) 上满足下式的线性形式，

\[
\mathrm{H} ^ {\prime} = \left\{x \in \mathrm{F} \mid \langle \alpha_ {\mathrm{H}}, x \rangle = 1 \right\},
\]

\(\mathfrak{H}_{\mathrm{H}}\) 就是超平面 \(\mathrm{L}_{\alpha_{\mathrm{H}},k} = \{x\in \mathrm{F}\mid \langle \alpha_{\mathrm{H}},x\rangle = k\}\)（其中 \(k\in \mathbf{Z}\)）所成的集合，而 \(\alpha_{\mathrm{H}}\) 与 \(-\alpha_{\mathrm{H}}\) 是具有此性质的仅有的线性形式。

因此，只要证明由形如 \(\pm\alpha_{H}\) 的 V 中元素组成的集合 R 是 V 中的约化根系，命题即得证。

\begin{enumerate}
\def\labelenumi{\alph{enumi})}
\item
  我们证明条件 \((\mathrm{RS}_{\mathrm{I}})\)：显然 R 是有限集（因为 \(H_{0}\) 是有限集）且不含 0。此外，R 生成 V。事实上，若 \(x \in F\) 与 R 正交，则 \(x \in H\)（对所有 \(H \in H_{0}\)），且沿向量 x 的平移与 G 的每个元素交换。由于 G 是本质的，这蕴含 x = 0。
\item
  我们证明 \((\mathrm{RS}_{\mathrm{II}})\)。对于 \(v \in V\) 和 \(r \in R\)，如上置 \(L_{v,r} = \{x \in F | \langle v, x \rangle = r\}\)；若 \(\alpha \in R\)，置 \(H_{\alpha} = L_{\alpha,0}\)，并令 \(s_{\alpha}\) 为 \(s_{H_{\alpha}}\) 的转置。存在唯一的 \(\alpha' \in F\) 与 \(H_{\alpha}\) 正交且满足 \(\langle \alpha', \alpha \rangle = 2\)。于是 \(s_{H_{\alpha}} = s_{\alpha', \alpha}\) 且 \(s_{\alpha} = s_{\alpha, \alpha'}\)。对于 \(\beta \in R\)，
\end{enumerate}

\[
\mathrm{L} _ {s _ {\alpha} (\beta), 1} = s _ {\mathrm{H} _ {\alpha}} \left(\mathrm{L} _ {\beta , 1}\right) \in \mathfrak {H}
\]

且存在 \(\gamma \in \mathbb{R}\) 和 \(n\in \mathbf{N}^*\)，使得 \(\mathrm{L}_{s_{\alpha}(\beta),1} = \mathrm{L}_{\gamma ,n}\)。于是

\[
s _ {\mathrm{H} _ {\alpha}} (\mathrm{L} _ {\gamma , 1}) = \mathrm{L} _ {\beta , 1 / n}
\]

所以 \(1 / n\in \mathbf{Z}\)。于是 \(n = 1\) 且 \(s_{\alpha}(\beta) = \gamma \in \mathbb{R}\)。这证明了 \((\mathrm{RS}_{\mathrm{II}})\)。

\begin{enumerate}
\def\labelenumi{\alph{enumi})}
\setcounter{enumi}{2}
\tightlist
\item
  我们证明 \((\mathrm{RS}_{\mathrm{III}})\)。令 \(\alpha \in R\)，并置 \(H'_{\alpha} = L_{\alpha,1}\)。于是
\end{enumerate}

\[
\mathrm{H} _ {\alpha} ^ {\prime} = \mathrm{H} _ {\alpha} + (1 / 2) \alpha^ {\prime};
\]

由于平移 \(t(\alpha^{\prime})\)（沿向量 \(\alpha^{\prime}\)）是乘积 \(s_{\mathrm{H}_{\alpha}^{\prime}}s_{\mathrm{H}_{\alpha}}\)（第五章第 2 节第 4 号命题 5），它属于 T，且 \(\alpha^{\prime} = t(\alpha^{\prime})(0)\) 是 G 的特殊点。因此，对所有 \(\beta \in \mathbb{R}\)，存在超平面 \(\mathrm{L}_{\beta,k}\) 经过 \(\alpha^{\prime}\)，其中 \(k\) 为整数；这说明 \(\langle \beta ,\alpha^{\prime}\rangle \in \mathbf{Z}\)，并证明了 (RSIII)。

\begin{enumerate}
\def\labelenumi{\alph{enumi})}
\setcounter{enumi}{3}
\tightlist
\item
  最后，显然 R 是约化的，因为若 \(H, H' \in H_{0}, H \neq H'\)，则线性形式 \(\alpha_{H}\) 与 \(\alpha_{H'}\) 不成比例。
\end{enumerate}

备注。1）当 G 不可约且无限时，T 的秩为 l 的假设特别得到满足。事实上，由 T 中平移所对应的向量生成的向量空间在 G 到 F 的线性群中的典范像下不变。当 G 无限时它不等于 \(\{0\}\)，故当 G 无限且不可约时它等于整个 F。

由反射生成的有限群不一定总是某个根系的韦伊群。更精确地说：

命题 9. 设 V 是 l 维实向量空间，G 是 \(\mathbf{GL}(V)\) 的一个有限子群，由反射生成且为本质的。在 V 上取一个 G 不变的标量积。以下条件等价：

\begin{enumerate}
\def\labelenumi{(\roman{enumi})}
\item
  存在 \(\mathrm{V}\) 的一个秩为 \(l\) 且在 \(\mathrm{G}\) 下不变的离散子群。
\item
  存在 V 上的一个 \(\mathbf{Q}\)-结构（《代数》第二章第 8 节第 1 号定义 1），且在 \(\mathbf{G}\) 下不变。
\item
  V 中存在一个以 G 为韦伊群的根系。
\item
  存在一个 V 的离散位移群 \(G'\)，它在 V 上适当作用并由反射生成，且 \(G'\) 是 G 与一个秩为 l 的平移群的半直积。
\item
  \(\Longrightarrow\) (i)：设 \(V' \subseteq V\) 是 V 上在 G 下不变的 Q-结构。令 A 为 \(V'\) 的有限子集，它生成 Q-向量空间 \(V'\)。将 A 替换为 \(\bigcup_{s \in G} s(A)\) 后，可设 A 在 G 下不变。令 B 为 V 中由 A 生成的子群。则 B 在 G 下不变、有限生成且无挠，因而存在一个 Z-基，它既是 \(V'\) 在 Q 上的一组基，也是 V 在 R 上的一组基。
\item
  \(\Longrightarrow\) (ii)：例如，这由第 1 节第 1 号命题 1 得出。
\item
  \(\Longrightarrow\) (iii)：令 \(G'\) 为满足条件 (iv) 的群。\(G'\) 的平移群的秩为 l，并且根据第五章第 3 节第 10 号命题 9，0 是 \(G'\) 的特殊点。命题 6 表明，存在约化根系 \(R_{0}\)，位于 \(V^{*}\) 中，使得 \(G'\) 可与 \(W_{a}(R_{0})\) 相识别；于是群 G 是 \(R_{0}\) 的对偶根系的韦伊群。
\item
  \(\Longrightarrow\) (iv)：假设 G 保持 V 的一个秩为 l 的离散子群 M 不变。对于任意反射 \(s \in G\)，有 \(s(x) - x \in M\)（对所有 \(x \in M\)），故直线 \(D_{s}\) 与 \(H_{s}\) 正交且与 M 相交；令 \(\alpha_{s}, -\alpha_{s}\) 为循环群 \(D_{s} \cap M\) 的生成元。由 \(\alpha_{s}\) 和 \(-\alpha_{s}\) 组成的集合 A 在 G 下不变，因而生成 M 的一个在 G 下不变的子群 \(M'\)；由于 G 是本质的，离散群 \(M'\) 的秩为 l。令 \(G'\) 为 V 的仿射变换群，它是 G 与向量属于 \(M'\) 的平移群的半直积。令 \(G_{1}'\) 为由 \(G'\) 的反射生成的 \(G'\) 的子群。我们将证明 \(G_{1}' = G'\)，这就完成了证明。首先，由于 G 由反射生成，\(G_{1}' \supseteq G\)。另一方面，对于 G 的任一反射 s，令 \(t_{s}\) 为沿向量 \(\alpha_{s}\) 的平移。变换 \(s \circ t_{s}\) 是一个反射，且 \(s \circ t_{s} \in G'\)；因此 \(t_{s}\) 是 \(G'\) 中两个反射的乘积；对 G 的每个反射 s 都如此，故向量属于 \(M'\) 的平移全都属于 \(G_{1}'\)。
\end{enumerate}

定义 2. 满足命题 9 的等价条件的群 G 称为晶体群。

备注。2）设 G 是由反射生成且本质的有限群。则 G 是晶体群，当且仅当其 Coxeter 矩阵的每个元素都是整数 1、2、3、4、6 中的一个。事实上，根据第 1 节第 5 号备注 3），该条件是必要的。它是充分的这一事实将由第 4 节给出的有限 Coxeter 群分类推出（直接证明见第五章第 4 节习题 6）。

